% SBES 2027 - Research Track (draft v0.1)
% DOUBLE-ANONYMOUS SUBMISSION: use \documentclass[sigconf,anonymous]{acmart}
% Camera-ready: remove "anonymous" and add \author block.
% Template: SBC/CBSoft adaptation of ACM acmart (NOT plain ACM_SigConf).
\documentclass[sigconf,anonymous]{acmart}

\setcopyright{none}
\acmConference[SBES 2027]{41st Brazilian Symposium on Software Engineering}{Month 2027}{City, Brazil}% <-- update when official call is out
\acmYear{2027}
\acmDOI{}

\title{Empirical Evaluation of Quality and Defects in Software Generated by AI Agents: A Comparative Study under Controlled Requirements}

\begin{abstract}
Source code generated by large language models (LLMs) has become routine in the software engineering (SE) lifecycle. The literature shows, in isolated fragments, that this code exhibits characteristic defect patterns, frequent security vulnerabilities, technical debt, tests with limited effectiveness, and imperfect reproducibility. No reviewed study, however, combines these dimensions in a single controlled chain: generation of \emph{complete systems} by different \emph{AI agents} from the \emph{same requirement}, under \emph{multiple executions}, evaluated by an \emph{independent oracle} covering functionality, defects, security, maintainability, and reproducibility. This paper presents a controlled, prospective, reproducible experiment filling that gap: four agents (distinct models on the same agent framework, 2$\times$ Nemotron, Ling, Mimo) each produced three independent executions of a complete system (task-management REST API, Python/FastAPI/PostgreSQL), totaling 12 deliveries evaluated by an independent 81-test oracle suite plus non-functional checks and a Tambon-taxonomy defect matrix with severity and inter-rater agreement. Results: only 3/12 deliveries booted; the 12 Tambon defects were 100\% Blocker/Critical; the agents' own tests detected \textbf{0/12} defects whereas the oracle detected \textbf{4/4} reachable (RQ6/H3); recurring patterns (missing dependency; \emph{passlib}+\emph{bcrypt} conflict; migration misconfiguration; ``ghost'' version) repeat across ``free'' models, suggesting a common ``quality trough''. The study descriptively supports that self-approval by an agent's own tests is no substitute for independent validation.
\end{abstract}

\keywords{AI-generated software; coding agents; software quality; software defects; independent oracle; software testing; controlled experiment}

\begin{document}

\maketitle

\section{Introduction}
LLMs now permeate the SE lifecycle, from code and test generation to program repair~\cite{A01}. Reviews consolidate recurring limitations---hallucinations, security weaknesses, generalization issues, interpretability~\cite{A02}---and the field keeps a conflict of results: favorable evidence (e.g., the GitHub Copilot controlled experiment~\cite{D03}) versus unfavorable (distinct defect profiles at scale~\cite{B03}; frequent vulnerabilities~\cite{E01}). The absence of consensus justifies a controlled empirical investigation.

Empirical studies show LLM-generated code has characteristic, recurring defects: Tambon et al.~\cite{B01} analyzed 333 bugs and proposed a ten-pattern taxonomy; Cotroneo et al.~\cite{B03} compared 500k+ human/AI samples and found \emph{different defect profiles}; LLM outputs are \emph{non-deterministic}~\cite{B04}. On testing, LLM-generated suites contain \emph{test smells}~\cite{C03}, have limited detection (29--60\% in Defects4J)~\cite{C04}, and coverage/mutation lose reliability when bugs exist~\cite{C05}; special values (\emph{None, inf, NaN}) are often omitted~\cite{C06}. On quality, non-functional aspects are underrated~\cite{D04} and AI code accumulates \emph{code smells}~\cite{D05}. On security, 24--30\% of snippets are vulnerable~\cite{E01}, all analyzed LLMs generate vulnerabilities~\cite{E02}, and security prompting does not reduce density~\cite{E03}. On coding agents, reproducibility is limited (68.3\% run in clean environments; 13.5$\times$ dependency expansion~\cite{F02}) and agents fail unevenly (commit reversion 0.7--7.6\%)~\cite{F05}.

\textbf{Gap.} None of the 25 reviewed studies combines, in one controlled and reproducible experiment, the same specification $\rightarrow$ different agents $\rightarrow$ complete systems $\rightarrow$ $\ge$3 executions $\rightarrow$ independent oracle $\rightarrow$ defects $+$ security $+$ maintainability $+$ reproducibility $\rightarrow$ statistical analysis. Retrospective studies observe real data; most work with snippets or patches, not complete systems; and few separate who \emph{generates} from who \emph{evaluates}.

\textbf{Research problem.} Stated with neutrality: although the literature shows that AI-generated code has characteristic defects and tests with low detection ability, there is no controlled empirical evidence on how different AI agents compare when producing \emph{complete systems} from the \emph{same requirements}, evaluated by an \emph{independent oracle}, under \emph{multiple executions}, covering functional quality, defects, security, maintainability, and reproducibility.

\textbf{Contributions.} (i)~A first controlled, prospective, reproducible experiment combining, in the same chain, complete systems $+$ different agents $+$ multiple executions $+$ independent oracle $+$ defects/NFR/testability; (ii)~evidence of \emph{recurring defect patterns transferable across ``free'' models} of the same agent pipeline (``quality trough''); (iii)~evidence on \emph{testability}: the agents' own tests detected no defect (0/12), while the independent oracle detected all reachable ones (4/4).

\section{Related Work}
Organized by theme (full entries with DOIs in References):

\begin{itemize}
\item \textbf{LLMs in SE.}~\cite{A01} largest SLR (395 articles);~\cite{A02} PRISMA review listing recurring limitations. They supply the RQ vocabulary but do not run controlled experiments on complete-system generation.
\item \textbf{Defects.}~\cite{B01} the adopted defect taxonomy.~\cite{B02} self-correction helps ($+$29.2\%)---a variable kept \emph{out} to preserve control.~\cite{B03} different profiles AI/human (motivates profile analysis).~\cite{B04} non-determinism (motivates $\ge$3 executions).
\item \textbf{LLM-based testing.}~\cite{C01} map of LLM testing.~\cite{C02} code+tests for small algorithms (this work advances to complete systems).~\cite{C03--C06} test smells, limited detection 29--60\%, unreliable coverage/mutation in the presence of bugs, omission of special values---support distrusting the AI's own tests.
\item \textbf{Quality \& security.}~\cite{D01,D02} supervision and context-dependence (neutrality).~\cite{D03} Copilot experiment.~\cite{D04} ISO/IEC 25010 adoption.~\cite{D05} 89.1\% \emph{code smells}.~\cite{E01--E03} vulnerability frequency and limited prompting gains.
\item \textbf{Coding agents.}~\cite{F01} agent model separation.~\cite{F02} reproducibility/dependency gap.~\cite{F03} requirement vs.\ intrinsic error.~\cite{F04,F05} uneven real behavior.
\end{itemize}

\textbf{Position.} The review (25 sources) shows 8\% generate a complete system, 12\% use an independent oracle, and 0\% combine complete systems $+$ defects $+$ security (gaps G1--G5). This experiment is positioned precisely in that gap.

\section{Methodology}
\subsection{Position, neutrality, pre-registration}
Controlled, prospective, reproducible experiment. Hypotheses, RQs, analysis plan, and the neutrality principle were fixed \emph{before} data collection and consolidated as a pre-registration with temporal anchors in the repository history. \textbf{Neutrality (C1):} the experiment does not start from ``AI is bad''; if data shows AI equal or better on a dimension, that is a result.

\subsection{Experimental design}
Between-groups comparison (one group per agent) with intra-group repetitions. \textbf{Experimental unit}: a complete generated system. \textbf{Treatment}: model identity inside the agent. \textbf{Replications}: 3 independent executions per agent (non-determinism~\cite{B04}). \textbf{Blinding}: oracle written by an independent QA team in a private, immutable repository, out of the agents' reach.

\textbf{Agents} (same framework \texttt{opencode}, varying only the model):

\begin{table}[t]
\caption{Agents (treatments).}
\label{tab:agents}
\small
\begin{tabular}{lll}
\toprule
\textbf{Treatment} & \textbf{Model} & \textbf{\texttt{opencode} ID}\\
\midrule
U & Nemotron 3 Ultra & \texttt{nemotron-3-ultra-free}\\
L & Nemotron 3.5 Lightning & \texttt{nemotron-3.5-lightning-free}\\
G & Ling 3.0 Flash & \texttt{ling-3.0-flash-fin-free}\\
M & Mimo v2.5 & \texttt{mimo-v2.5-free}\\
\bottomrule
\end{tabular}
\end{table}

Model versions are recorded per execution in manifests (agents have versions~\cite{F05}). Limitation: agent-tool comparison is lost (mitigated by the identical loop).

\subsection{Specification (the ``same requirement'')}
Target system: \emph{task-management REST API} (moderate scope), Python + FastAPI + PostgreSQL. FR1--FR8: registration/login (hash + JWT); task CRUD; permission rules (owner/admin); filters/pagination; input validation; standardized errors; PostgreSQL with migrations; automated tests in the delivery. NFRs anchored in ISO/IEC 25010~\cite{D04} and CWE/OWASP: security, maintainability, reproducibility, basic performance---informed only with behavioral acceptance criteria, \emph{without} solution hints, to avoid inducing groups unequally. This document is the only input to the agents.

\subsection{Independent oracle}
The \textbf{81-test} suite is independent of the generating AI, private and immutable during collection, validated against a golden (non-AI) reference. Layers: functional (pytest covering FR1--FR8); negative (invalid/extreme inputs, permissions); security (bandit, pip-audit); structural (Ruff + radon); testability (agent's tests run against the oracle). Coverage is \emph{reported but not treated as quality evidence}~\cite{C05}; real defect detection counts. \textbf{Environment (D9):} validation ran \emph{native} (VM without nested virtualization; PostgreSQL + venv + \texttt{alembic upgrade head} + \texttt{/health} gate); Docker artifacts remain required and are validated structurally, to be executed later on the production machine---reported as a limitation (Sec.~5.2).

\subsection{Defect matrix}
Each defect: ID; agent/execution; location; \emph{category} (Tambon, 10 patterns~\cite{B01}); \emph{severity} (Blocker/Critical/Major/Minor, harmonized with CWE); \emph{detection source} (oracle / agent's tests); reproduction evidence. \textbf{Independent double classification} (two raters, Cohen's $\kappa$, divergence resolution recorded). Density = defects per KLOC.

\subsection{Metrics and statistical analysis}
RQ/H $\rightarrow$ metric $\rightarrow$ instrument: RQ1 pass rate (pytest); RQ2/H1 defects/KLOC (descriptive + Poisson CI); RQ3/H2 category distribution ($\chi^2$ + Cram\'er's V; pointwise Fisher); RQ4 severity (descriptive); RQ5/H4 functional vs.\ NFR (per-dimension, no single index); RQ6/H3 detection matrix ($2\times2$, Fisher exact); RQ7 reproducibility (executability check); RQ8 variability (CV). $\alpha=0.05$ with multiple-comparison correction; $\chi^2$ only conclusive with expected cells $\ge5$, otherwise Fisher exact \emph{reported as informative, not conclusive}. Negative results are results (neutrality).

\section{Results}
Twelve deliveries (4 agents $\times$ 3 executions); calibration pilots excluded. Figures~1--5 are generated by a reproducible script from the aggregates.

\subsection{Functional quality (RQ1)}
\begin{table}[t]
\caption{Oracle execution per delivery.}
\label{tab:rq1}
\small
\begin{tabular}{llllll}
\toprule
\textbf{Delivery} & \textbf{Tests} & \textbf{Fail.} & \textbf{Err.} & \textbf{Pass} & \textbf{Boot}\\
\midrule
lightning/e1--e2 & 0 & 0 & 0 & 0 & no\\
\textbf{lightning/e3} & 81 & 32 & 48 & \textbf{1} & \textbf{yes}\\
ling/e1--e3 & 0 & 0 & 0 & 0 & no\\
\textbf{mimo/e2--e3} & 81 & 21 & 48 & \textbf{12} & \textbf{yes}\\
ultra/e1--e3 & 0 & 0 & 0 & 0 & no\\
\bottomrule
\end{tabular}
\end{table}

Only \textbf{3/12 booted}; pass rates 1/81 and 12/81---no delivery passed 100\% (Fig.~\ref{fig:boot}).

\begin{figure}[t]
\includegraphics[width=\columnwidth]{figuras/fig1_densidade}
\caption{Defect density per agent (defects/KLOC).}
\label{fig:density}
\end{figure}
\begin{figure}[t]
\includegraphics[width=\columnwidth]{figuras/fig2_heatmap_categoria_agente}
\caption{Heatmap of defect category $\times$ agent.}
\label{fig:heatmap}
\end{figure}
\begin{figure}[t]
\includegraphics[width=\columnwidth]{figuras/fig3_bootabilidade}
\caption{Bootability and oracle pass rate.}
\label{fig:boot}
\end{figure}
\begin{figure}[t]
\includegraphics[width=\columnwidth]{figuras/fig4_kappa}
\caption{Inter-rater agreement ($\kappa$).}
\label{fig:kappa}
\end{figure}
\begin{figure}[t]
\includegraphics[width=\columnwidth]{figuras/fig5_severidade}
\caption{Severity distribution.}
\label{fig:severity}
\end{figure}

\subsection{Defects and density (RQ2/H1; RQ3/H2)}
\textbf{12 Tambon defects} (2 NFR/non-bug and 2 infrastructure findings excluded). Density: ultra 4/4.30=0.93; lightning 4/2.93=1.36; ling 1/1.83=0.55; mimo 3/3.27=0.92 (Fig.~\ref{fig:density}). With $\le$4 defects per agent, H1 is \emph{descriptive} (wide Poisson CI).

\begin{table}[t]
\caption{Defect category profile per agent.}
\label{tab:h2}
\small
\begin{tabular}{lcccccc}
\toprule
\textbf{Agent} & \textbf{inc-gen} & \textbf{silly} & \textbf{w-input} & \textbf{nprompt} & \textbf{haluc} & \textbf{w-attr}\\
\midrule
ultra & 4 & 0 & 0 & 0 & 0 & 0\\
lightning & 0 & 1 & 0 & 1 & 1 & 1\\
ling & 0 & 0 & 1 & 0 & 0 & 0\\
mimo & 1 & 2 & 0 & 0 & 0 & 0\\
\bottomrule
\end{tabular}
\end{table}

$\chi^2$ agent$\times$category: \textbf{$\chi^2=25.7$, df=15, $p=0.041$, V=0.85}---informative, not conclusive (cells $<5$). Ultra concentrated 100\% in \texttt{incomplete\_generation} (pointwise Fisher \textbf{$p=0.010$}), the most sustainable finding (Fig.~\ref{fig:heatmap}).

\subsection{Severity (RQ4)}
All 12 defects were \textbf{Blocker (8) or Critical (4)} (Fig.~\ref{fig:severity}). This is expected and inflated by the detection bias: 10/12 deliveries did not boot.

\subsection{Recurring patterns}
\begin{table}[t]
\caption{Recurring defect patterns across the 12 deliveries.}
\label{tab:patterns}
\small
\begin{tabular}{lll}
\toprule
\textbf{Pattern} & \textbf{Occurrences} & \textbf{Examples}\\
\midrule
Missing dependency (\texttt{pydantic-settings}/\texttt{email-validator}) & 4 (+1 pilot) & ultra e1/e3, mimo e1\\
\emph{passlib} 1.7.4 $+$ \emph{bcrypt} $\ge$4.4 $\rightarrow$ 500 on \texttt{/auth/register} & 3 (+1 pilot) & lightning e3, mimo e2/e3\\
Migration misconfig. (alembic w/o \texttt{script\_location}) & 1 (+1 pilot) & ultra e2\\
``Ghost'' version (\texttt{pydantic==2.8.4} non-existent) & 1 & lightning e1\\
\bottomrule
\end{tabular}
\end{table}

These patterns \textbf{repeat across distinct ``free'' models} of the same pipeline---a common ``quality trough'' at collection time, regardless of model family.

\subsection{Testability --- can the agent's own tests detect the defects? (RQ6/H3)}
Per-defect $2\times2$ matrix (formal executions):

\begin{table}[t]
\caption{Detection: agent's own tests vs.\ independent oracle.}
\label{tab:h3}
\small
\begin{tabular}{lccc}
\toprule
 & \textbf{Agent tests: yes} & \textbf{Agent tests: no} & \textbf{Total}\\
\midrule
\textbf{Oracle: detected} & 0 & \textbf{4} & 4\\
\textbf{Oracle: not detected} (non-bootable) & 0 & 8 & 8\\
\textbf{Total} & \textbf{0} & \textbf{12} & 12\\
\bottomrule
\end{tabular}
\end{table}

The agents' own tests detected \textbf{0/12} defects (0/18 matrix lines, including NFR and pilot). In the function-exercisable deliveries, the oracle detected \textbf{4/4} reachable defects (lightning/e3 $\times$2; mimo/e2, e3) while the agent's tests detected 0---even when the agent suite could run against the oracle. The extreme marginal (0-vs-4) is reported as a \emph{strong trend, not significance} (n, selection bias). The transversal datum: \emph{approval by an agent's own tests was, in no case, a guarantee of conformance with the oracle}, consistent with~\cite{C03,C04,C05,C06}.

\subsection{Inter-rater agreement and reproducibility (RQ7--RQ8)}
Over 18 matrix items (blinded double classification): \textbf{$\kappa$ category = 0.54 (moderate)}; \textbf{$\kappa$ severity = 0.91 (almost perfect)} (Fig.~\ref{fig:kappa}). The 6 category divergences concentrate on dependency-integration defects (resolved in a recorded meeting). Reproducibility (RQ7): 3/12 ran in a clean environment; the rest failed on the Sec.~4.4 patterns. Variability (RQ8): mimo stable across two runs (pass rate 0.15/0.15); lightning single run; ultra/ling no functional signal.

\subsection{Global quality: functional $\ne$ non-functional (RQ5/H4)}
Bootable deliveries still fail 69--80/81 tests (auth/DB defects); non-bootable ones may satisfy static NFRs. A single index would collapse distinct dimensions---descriptively supporting H4.

\section{Discussion}
\subsection{Interpretation}
\textbf{Profile, not quantity.} Density differences are descriptive; the sustainable finding is ultra's 100\% \texttt{incomplete\_generation} profile ($p=0.010$), coherent with evidence that AI and humans---and here distinct models---have distinct defect profiles~\cite{B03}. \textbf{H1 neither rejected nor accepted} with this sample; \textbf{H2} finds pointwise support.

\textbf{Testability (RQ6/H3).} The central result: \textbf{0/12} by the agents' own tests vs.\ \textbf{4/4} by the oracle. Self-approval $\ne$ independent validation, in line with the literature~\cite{C03,C04,C05,C06}. The asymmetry (8/12 without functional oracle execution) limits the comparison.

\textbf{Functional $\ne$ global (RQ5/H4).} Dimensions diverge; a single index would collapse them~\cite{D04}.

\textbf{Quality trough.} Recurring patterns extend, to the controlled ``free''-model scenario, the findings of limited reproducibility/executability~\cite{F02} and unequal agent reliability~\cite{F05}.

\subsection{Central limitations (honest reading)}
\begin{enumerate}
\item \textbf{Statistical power} ($n=12$; expected cells $<5$): H1--H3 are trends, not conclusions.
\item \textbf{Functional evaluation on 3/12}: latent defects in non-bootable deliveries were unmeasurable.
\item \textbf{Environment conditioning} (native validation, no Docker): container-only configurations were treated as portability defects, separating ``defect'' from ``environment assumption'' (D9).
\item \textbf{Manual classification}: moderate category $\kappa$ (0.54).
\item \textbf{``Free'' models at collection time}: results describe these models on that date; versions/routes change the picture.
\item \textbf{Blocker bias} from non-bootability.
\item \textbf{No generalization}: one system/stack (tasks + Python/FastAPI/PostgreSQL).
\end{enumerate}

\subsection{Synthesis}
In the controlled scenario, free-model agents produced complete systems with low execution rate (3/12), defects concentrated in a few recurring patterns, and no full conformance with the oracle. The results support neither ``AI does not work'' nor ``AI replaces QA'': they support that, under controlled requirements and an independent oracle, the quality delivered by these agents at collection time \emph{fell below the functional minimum in 9 of 12 executions}, and that approval by the agents' own tests is no substitute for independent validation.

\section{Conclusion and Future Work}
This study contributes a controlled, prospective, reproducible experiment combining, for the first time as far as the review reaches, complete-system generation by different AI agents from the same requirement, multiple executions, an independent oracle, and simultaneous functional, defect, security, maintainability, reproducibility, and testability evaluation. Highlights: low executability (3/12); a defect profile concentrated in recurring patterns (``quality trough''); and \emph{no defect detected by the agents' own tests (0/12)} vs.\ oracle detection of all reachable (4/4). Neutrality upheld: the data, not a premise, support these readings, within the power/generalization limits.

\textbf{Future work.} (i)~Expand the sample (more executions; Docker execution on the production machine) to reach expected cells $\ge5$; (ii)~run the delivered containers; (iii)~extend to more models/families and to distinct agents (comparing the tool, not only the model); (iv)~full pre-registration of the expanded design; (v)~longitudinal check of the ``quality trough'' across model updates; (vi)~minimal human-intervention contrast; (vii)~pre-registered laboratories for the type/domain and boundary hypotheses; (viii)~a full formal SLR.

\section*{Availability of Artifacts}
The aggregated data used in Figures~1--5, the figure-generation script, and summarized results are available to reviewers in an \emph{anonymized repository} (anonymous.4open.science); instructions to reproduce the tables are in the package. The executable specification and the private oracle suite are intentionally \emph{not} public at review time to preserve double-anonymous review and avoid exposing internal test artifacts; after acceptance, aggregated, curated artifacts will be deposited with a DOI under CC-BY-4.0, and raw execution logs (containing environment/app credentials) remain restricted with documented justification. License: CC BY 4.0.

\begin{thebibliography}{99}
\scriptsize

\bibitem{A01} HOU, X. et al.\ \emph{Large Language Models for Software Engineering: A Systematic Literature Review}. ACM ToSEM, v.~33, n.~8, art.~220, 2024. DOI: 10.1145/3695988.
\bibitem{A02} UMAMA et al.\ \emph{LLM-Based Code Generation: A Systematic Literature Review With Technical and Demographic Insights}. IEEE Access, v.~13, pp.~194915--194939, 2025. DOI: 10.1109/ACCESS.2025.3631952.
\bibitem{B01} TAMBON, F. et al.\ \emph{Bugs in Large Language Models Generated Code: An Empirical Study}. Empirical Software Engineering, v.~30, art.~65, 2025. DOI: 10.1007/s10664-025-10614-4.
\bibitem{B02} DOU, S. et al.\ \emph{What's Wrong with Your Code Generated by Large Language Models? An Extensive Study}. arXiv:2407.06153, 2024.
\bibitem{B03} COTRONEO, D.; IMPROTA, C.; LIGUORI, P. \emph{Human-Written vs.\ AI-Generated Code: A Large-Scale Study of Defects, Vulnerabilities, and Complexity}. IEEE ISSRE 2025, pp.~252--263. DOI: 10.1109/ISSRE66568.2025.00035.
\bibitem{B04} OUYANG, S.; ZHANG, J. M.; HARMAN, M.; WANG, M. \emph{An Empirical Study of the Non-Determinism of ChatGPT in Code Generation}. ACM ToSEM, v.~34, n.~2, art.~42, 2025. DOI: 10.1145/3697010.
\bibitem{C01} WANG, J. et al.\ \emph{Software Testing With Large Language Models: Survey, Landscape, and Vision}. IEEE TSE, v.~50, n.~4, pp.~911--936, 2024. DOI: 10.1109/TSE.2024.3368208.
\bibitem{C02} BEER, A. et al.\ \emph{Examination of Code generated by Large Language Models}. arXiv:2408.16601, 2024.
\bibitem{C03} ALVES, V. A.; SANTOS, C.; BEZERRA, C. I. M.; MACHADO, I. \emph{Detecting Test Smells in Python Test Code Generated by LLM: An Empirical Study with GitHub Copilot}. SBES 2024. DOI: 10.5753/sbes.2024.3561.
\bibitem{C04} YANG, L. et al.\ \emph{On the Evaluation of Large Language Models in Unit Test Generation}. ASE 2024. DOI: 10.1145/3691620.3695529.
\bibitem{C05} ZHAO, J.; ZHOU, S.; COHEN, E. \emph{Do Coverage and Mutation Scores of LLM-Generated Test Suites Correlate with Their Effectiveness? (Replicability Study)}. ISSTA 2026. DOI: 10.1145/3832093.
\bibitem{C06} WALCZAK, J.; TOMALAK, P.; LASKOWSKI, A. \emph{Impact of code context and prompting strategies on automated unit test generation with modern general-purpose large language models}. JSS, v.~237, art.~112834, 2026. DOI: 10.1016/j.jss.2026.112834.
\bibitem{D01} TOSI, D. \emph{Studying the Quality of Source Code Generated by Different AI Generative Engines: An Empirical Evaluation}. Future Internet, v.~16, n.~6, art.~188, 2024. DOI: 10.3390/fi16060188.
\bibitem{D02} MOLISON, A. S. et al.\ \emph{Is LLM-Generated Code More Maintainable \& Reliable than Human-Written Code?}. arXiv:2508.00700, 2025.
\bibitem{D03} BAUER, J.; LINDEMAN, L.; REILLY, L.; RODRIGUEZ, M. \emph{Does GitHub Copilot improve code quality? Here's what the data says}. GitHub Customer Research, 2024.
\bibitem{D04} SUN, X.; ST\AA{}HL, D.; SANDAHL, K.; KESSLER, C. \emph{Quality assurance of LLM-generated code: Addressing non-functional quality characteristics}. JSS, v.~238, art.~112885, 2026. DOI: 10.1016/j.jss.2026.112885.
\bibitem{D05} LIU, Y. et al.\ \emph{Debt Behind the AI Boom: A Large-Scale Empirical Study of AI-Generated Code in the Wild}. arXiv:2603.28592, 2026.
\bibitem{E01} FU, Y. et al.\ \emph{Security Weaknesses of Copilot-Generated Code in GitHub Projects: An Empirical Study}. ACM ToSEM, v.~34, n.~8, 2025. DOI: 10.1145/3716848.
\bibitem{E02} MORKONDA, S. G.; SELIM, M.; ASSAL, H. \emph{Security of LLM-generated Code: A Comparative Analysis}. arXiv:2605.23091, 2026.
\bibitem{E03} KHARMA, M. et al.\ \emph{An Empirical Evaluation of LLM-Generated Code Security Across Prompting Methods}. arXiv:2605.24298, 2026.
\bibitem{F01} HORIKAWA, K. et al.\ \emph{Agentic Refactoring: An Empirical Study of AI Coding Agents}. arXiv:2511.04824, 2025.
\bibitem{F02} VANGALA, B. P.; ADIBIFAR, A.; GEHANI, A.; MALIK, T. \emph{AI-Generated Code Is Not Reproducible (Yet): An Empirical Study of Dependency Gaps in LLM-Based Coding Agents}. arXiv:2512.22387, 2025.
\bibitem{F03} BELOZEROV, V.; BARCLAY, P. J.; SAMI, A. \emph{Secure coding with AI --- from detection to repair}. Empirical Software Engineering, v.~31, art.~93, 2026. DOI: 10.1007/s10664-026-10812-8.
\bibitem{F04} EHSANI, R. et al.\ \emph{Where Do AI Coding Agents Fail? An Empirical Study of Failed Agentic Pull Requests in GitHub}. MSR 2026, pp.~807--811. DOI: 10.1145/3793302.3793579.
\bibitem{F05} OUKHAY, I.; BEGOUG, M.; CHOUCHEN, M.; OUNI, A. \emph{When AI Code Doesn't Stick: An Empirical Study on Reverted Changes Introduced by AI Coding Agents}. MSR 2026, pp.~847--851. DOI: 10.1145/3793302.3793587.

\end{thebibliography}

\end{document}